Onder een vilten hoed
zat een geslepen snoet

vaardig beschermd tegen
de zon en de regen.

Ook gaf een wollen jas,
die flink versleten was,

nog genoeg bescherming
tegen de marteling

van het Noord-Duitse weer
voor de trekkende heer.

Eronder zat een vest
ook niet meer al te best.

Daaronder nog een hemd.
Zo werd de kou getemd.

De talrijke zakken
met legio vakken

zaten vol met koopwaar
en werden zo erg zwaar,

maar zijn kleding zat los.
Bovendien trok zijn ros,

door zijn kar te vullen,
het gros van de spullen.

Op zijn rug nog een mand
en voor zich aan de hand

nog zijn oude handkar.
Het was nogal bizar
dat hij die had bewaard,
nadat hij had gespaard

voor zijn paard en wagen.
Op zulke navragen

verried hij zijn hebzucht
verheerlijkt in zelftucht.

Dus moest hij veel lopen
om zat te verkopen.

Een kniebroek van linnen,
om maar te beginnen,

diende de wandelgang.
Sok of kous naar gelang

de omstandigheden.
Ook de schoenen deden

hun best de zware last
van de schriele marsgast

gedienstig te dragen
zonder maar te klagen

over het natte weer
dat door het lekke leer

naar binnen sijpelde.
Gierigheid twijfelde

niet aan een vervanging.
Waardoor de smart doorging.

Hij lichtte mensen op.
Hij strikte koud hun strop
door troep aan te smeren.
Het kon hem niet deren

als ze te arm waren
voor dure koopwaren.

Hij was een verkoper,
die een rode loper

kreeg uitgerold bij hen
die met hun kippenren

last hadden van vossen.
Hij kon dat oplossen

met een bepaalde stof.
Die vond een vos niet tof.

De stof was wel speciaal.
De prijs was abnormaal

om die valse reden.
De vossen vermeden

het hok natuurlijk niet.
Maar de deal was geschied.

Zo had hij meer leugens
om de domste wezens

geld af te troggelen.
Hij bleef sluw broddelen.

Emmy kon hem doorzien,
maar zij had geen aanzien

in haar boerengezin.
Blind door beloofd gewin
stonken ze steeds erin.
De marktman kreeg zijn zin

en vernauwde zijn fuik.
Hij maakte sluw misbruik

van de zonde hebzucht
en pleegde zelfs ontucht,

in een milde vorm dan,
slinks onder het mom van

een speelse affectie,
met zijn vijand Emmy

door kort in haar billen
zijn wellust te stillen.

Hij kneep haar bedreven.
Emmy stond te beven

van woede en walging
over zijn aanranding.

Plots stopte de venter.
Er kwam iets urgenter

dwars op zijn handelspad.
De marskramer zag dat

twee dode lichamen
in een greppel samen

met een lege rijkoets
lagen die onverhoeds

gans waren leeggeroofd.
Hij die in praal geloofd

zag een gelegenheid
voor zijn hebberigheid.

Hij vond in de restjes
nieuwe schoenen, vestjes,
wielen voor zijn wagen
of geld voor kon vragen.

Een fluwelen kussen
zat daar mede tussen

met “Adelheid Schaffgotsch”
als adellijke trots

er fijn in gegraveerd.
Die vrouw was hem gepeerd

zo vond de trekvent uit.
Hij vergaarde zijn buit

en vertrok zonder stijl,
lomp en in aller ijl,

de plaats van het misdrijf,
uit angst dat het kakwijf

nog terug zou keren
en hem zou beleren.

Kruis kopen dus doen of je gelovig bent is goedkoper dan tol
